\begin{definition}{(vollst\"andig) unabhängig/identisch verteilt/zentriert}
                  {UnabhaengigIdentischVerteilt}
Seien $n, m \in \N$, $\gamma \in \mcm_n(\Omega, \mfa, \R^n)$,
$\delta \in \mcm_m(\Omega, \mfa, \R^m)$ und $F: \R \rightarrow [0, 1]$ eine
Verteilungsfunktion.\\
$\gamma$ heißt \textbf{unkorreliert}
\[: \iff \gamma \in L_2(\Omega, \mfa, \Prim, \R^n) \; \wedge \;
    \forall\, i, j \in \haken n: i \neq j \folgt \E(\gamma_i\gamma_j)
  = \E(\gamma_i)\E(\gamma_j)\text.\]
$\gamma$ und $\delta$ heißen \textbf{vollständig unkorreliert}
\[:\iff \gamma, \delta \text{ unkorreliert } \; \wedge \;
        \forall\, i \in \haken n\, \forall\, j \in \haken m:
        \E(\gamma_i\delta_j)= \E(\gamma_i)\E(\delta_j)\text.\]
$\gamma$ heißt \textbf{unabhängig}
\[:\iff \forall A \in \mathfrak B^n: 
\Prim\left(\bigcap_{i=1}^n \meng{\gamma_i\in A_i}\right)
=\prod_{i=1}^n \Prim(\gamma_i\in A_i)\text.\]
$\gamma$ und $\delta$ heißen \textbf{vollständig unabhängig}
\begin{align*}
&:\iff \,
\forall\, A\in \mfb^n\, \forall\, B\in \mfb^m:\\
& \Prim\left(\bigcap_{i=1}^n \meng{\gamma_i \in A_i} \cap 
\bigcap_{j=1}^m \meng{\delta_j \in B_j} \right)
= \left(\prod_{i=1}^n \Prim(\gamma_i\in A_i)\right)
\left(\prod_{j=1}^m \Prim(\delta_j\in B_i)\right)\text.
\end{align*}
$\gamma$ heißt \textbf{identisch verteilt}
$: \iff \forall \, i, j \in \haken n\, \forall \, t \in \R:
    \Prim(\gamma_i \le t) = \Prim(\gamma_j \le t)$.\\
$\gamma$ heißt \textbf{$F$--verteilt}
$:\iff \gamma \text{ identisch verteilt } \; \wedge \;
 \forall \, t \in \R: \Prim(\gamma_1 \le t) = F(t)$.\\
$\gamma$ heißt \textbf{zentriert}
$:\iff \forall\, i \in \haken n: \E(\gamma_i) = 0$.\\
$\gamma$ heißt \textbf{$L_2$--normiert}
$:\iff \gamma \in L_2(\Omega, \mfa, \Prim, \R^n) \; \wedge \;
   \forall i \in \haken n: \norm{f_i}_2 = 1$.
\end{definition}